\chapter{Conclusion and Future Work}
\label{ch:conc}

\section{Conclusion}

This thesis presented a multimodal Transformer-oriented pipeline for analyzing post-pandemic vaccine-related misinformation dynamics.
On CONSTRAINT, a RoBERTa text classifier achieved macro-F1 $0.9719$.
On MMCoVaR, ablation showed that text is the dominant modality (macro-F1 $\approx 0.963$), image-only performance is weak, and late fusion approximately matches text without clear gains.
BERTopic analysis of 5{,}752 misinformation documents produced a multi-narrative taxonomy spanning sensational statistics, politicization, contested measures, origin conspiracies, miracle cures, and elite distrust---themes relevant to vaccine hesitancy communication.

Collectively, the work delivers a reproducible experimental stack and evidence-based answers to the research questions, while honestly scoping platform and modality limitations.

\section{Future work}

\begin{enumerate}
  \item Integrate TikTok video frames and Whisper audio transcripts for true short-video multimodality.
  \item Explore cross-attention / bottleneck fusion and unfrozen vision towers.
  \item Expand claim-level fact-checked labels beyond publisher reliability.
  \item Build a multilingual (including Vietnamese) evaluation set for USTH/local relevance.
  \item Deploy an interactive dashboard linking classifier scores to live narrative trends for public-health partners.
\end{enumerate}

\section{Final remarks}

Multimodality remains a promising direction for social misinformation research, but gains are dataset-dependent.
For text-heavy news, investing first in strong language models and narrative analytics may yield more actionable insight than complex fusion alone.
The pipeline and findings in this thesis provide a foundation for that continued work.
